\section{Discussion}\label{sec:discussion}

\paragraph{What is established.} Four statements are supported, with qualifications. First, adaptive conformal intervals were better calibrated than a calibration performed once and never updated; the direction is the same in the exploratory data, on the first hold-out panel at both warm-ups, for three base forecasters and after removing any single country, but R1 was met on that panel only at warm-up 96, a sensitivity run, and at warm-up 48 the intervals include zero (\ci{-5.4}{+0.5} and \ci{-3.9}{+1.2}). The statement is one of consistent direction, not of a pre-specified confirmation on that panel; on the second panel R1 was met at warm-up 48 for all four base models and not at 96. Second, in a secondary post hoc analysis ACI with ARIMA had a lower CE than the native ARIMA intervals at every horizon in both samples (the interval includes zero at $h=6$ and $h=12$ on the first hold-out panel). Third, at 12 months pooled coverage is below nominal for every method on the design and first hold-out panels, and delayed feedback explains part of the shortfall, by an amount that depends on method and sample. Fourth, ACI, DRACP and FACI have a median width below rolling at short horizons, but the pooled mean width of ACI at $h=3$ exceeds rolling's in both samples and the Winkler ratio is below 1 only on the hold-out panel at $h=1$ and $h=3$ (0.95 and 0.99 for ACI, the latter with an interval that includes zero); no efficiency advantage is established.

\paragraph{What is not established.} (i) Whether ACI improves on a rolling window depends on the panel: not shown on the design and first hold-out panels (point estimates of about 1 to 3~pp in the direction of an advantage, MDE undefined, inconclusive), not significant on the earlier 20 countries re-run from 2010-01 (Holm $p$ 0.22 to 1.0, MDE undefined, inconclusive under E5), and met at warm-up 48 on the second panel at $h=3$ for LightGBM, ARIMA and naive B and at $h=6$ for ARIMA and naive B, but inconclusive for Chronos at $h=3$ and not met for Chronos and LightGBM at $h=6$. The reason is unknown; the supported statement is that the result is panel dependent, and it is neither ``ACI does not improve on a rolling window'' nor ``ACI improves on a rolling window''. (ii) An advantage of DRACP over ACI was not detected and equivalence was not stable across warm-ups (R3); on the second panel DRACP had a lower CE in every cell at warm-up 48, by 0.5 to 1.8~pp, with no Holm $p$ below 0.05. (iii) That delay explains the $h=12$ shortfall was not met at the threshold (R4). None of these is evidence that the effects are absent; ``DRACP does not transfer'' and ``ACI and DRACP are equal'' are not allowed by the data.

\paragraph{Why null results.} The data do not apportion the reasons. The first hold-out panel has 7 countries and about 200 origin months at $h=3$, errors of $h$-step forecasts overlap, adaptive methods have a memory of about 50 months, and countries share shocks, so the block-bootstrap intervals are about 2 to 6~pp wide. There is a ceiling effect: the CE of rolling on that panel is already small (0.029, 0.036, 0.032 at $h=1,3,6$ against 0.020, 0.020, 0.025 for ACI), the power simulation cannot create an improvement larger than the baseline CE, and 80 percent power was not reached up to 10~pp, so we read the undefined MDE as a limit of the design and not as a finding about the methods. The 2~pp margin is about the size of ACI's own CE, so ``equivalent'' is weak by construction and the null outcomes against ACI were likely before the run; the informative contrasts are those against weak references, split (CE 0.050 to 0.056) and the native intervals. A true small difference cannot be separated from low power. For DRACP, regime and scale models are trained on half of the warm-up (24 rows at warm-up 48), its results depend on the warm-up, and its authors state that its advantage is specific to one-step-ahead forecasting, which our multi-month design does not test; our partial one-step reproduction could not distinguish ACI and DRACP either.

\paragraph{The panel dependence of R2.} The second panel is a counter-example to the null outcomes and is reported beside them. Candidate explanations are a shorter series, lower-income and more volatile economies, a warm-up occupying a larger share of 198 months, and a rolling window with a larger CE (0.066 to 0.090 at $h=3$ against 0.036 on the first panel; the MDE is defined here, 1.2 to 6.6~pp, which is consistent with the ceiling argument), but none was tested and warm-up 96 cannot arbitrate. The panels are not averaged and the favourable panel is not selected; a reader who wants a recommendation for a new set of countries cannot take it from this study.

\paragraph{Calibrated is not useful, and practice.} A nominal-coverage interval can be very wide: the pooled mean width of ACI is 4.28 pp at $h=3$ and 15.5 at $h=12$ in the exploratory data, and at $h=12$ an interval with honest 90 percent coverage may be too wide for practical use. A practitioner should not calibrate once and leave the calibration unchanged: one-shot calibration had a higher CE than ACI in every reported contrast for LightGBM, ARIMA and naive B and for Chronos in all but one cell (second panel, warm-up 96, $h=6$, $+0.8$~pp, indistinguishable from zero), and its coverage varied most across countries (standard deviation at $h=12$ 0.152 for split against 0.058 for ACI), although R1 was met on the first panel only at the sensitivity warm-up. Between a rolling window and ACI the study cannot recommend one on statistical or efficiency grounds. At 12 months users should treat the nominal level as an upper bound, since part of the shortfall comes from the delay in outcomes and cannot be removed in real use. The hold-out countries share calendar months with the design sample, so agreement between samples is weaker evidence of generality than if they had been observed in other periods.
